%% ============================================================
%% Degradation Mechanisms and Maintenance Strategies for
%% Composite Silicone Rubber Insulators in Electrified Railway
%% Catenary Systems: A Systematic Review
%%
%% Authors: Kusyama Lameck Kusyama, Frank Lujaji
%% Target: Engineering Failure Analysis (Elsevier) -- Review Article
%% Template: Elsevier elsarticle (numbered citations)
%% Compile on Overleaf or local TeX Live with pdflatex
%% ============================================================

\documentclass[preprint,12pt,3p]{elsarticle}

\usepackage{lineno,hyperref}
\usepackage{booktabs}
\usepackage{array}
\usepackage{graphicx}
\usepackage{amsmath}
\usepackage{longtable}
\usepackage[margin=2.5cm]{geometry}

\modulolinenumbers[5]

\journal{Engineering Failure Analysis}

\bibliographystyle{elsarticle-num}

%% ============================================================
\begin{document}

\begin{frontmatter}

\title{Degradation Mechanisms and Maintenance Strategies for Composite
Silicone Rubber Insulators in Electrified Railway Catenary Systems:
A Systematic Review}

\author[dit]{Kusyama Lameck Kusyama}
\author[dit]{Frank Lujaji\corref{cor1}}
\cortext[cor1]{Corresponding author. Tel.: +255-xxx-xxxxxx}
\ead{frank.lujaji@dit.ac.tz}

\address[dit]{Department of Mechanical and Industrial Engineering,
Dar es Salaam Institute of Technology, P.O. Box 2958, Dar es Salaam, Tanzania}

\begin{abstract}
Composite silicone rubber insulators (CSRIs) have displaced ceramic and
glass counterparts in electrified railway catenary systems over the past
three decades, principally on account of their hydrophobic surface
properties, reduced mass, and superior performance under polluted
conditions. Despite these advantages, service experience worldwide
confirms that CSRIs are susceptible to a range of degradation mechanisms
that progress at rates strongly dependent on the operating environment.
In tropical and coastal Sub-Saharan African contexts---conditions that
characterise emerging railway corridors such as Tanzania's Standard
Gauge Railway (SGR)---the relevant literature is sparse and no validated
maintenance management model tailored to such environments exists. This
paper presents a systematic review of the published literature on CSRI
degradation mechanisms, diagnostic techniques, and maintenance strategies,
with the objective of establishing the empirical foundation for a
context-specific maintenance framework applicable to tropical electrified
railway catenary systems. A structured search of four databases
(Scopus, Web of Science, IEEE Xplore, and Google Scholar) returned
1,247 records, of which 58 studies met the inclusion criteria after
two-stage screening. The review identifies five principal degradation
pathways: hydrophobicity loss, surface erosion and tracking, ultraviolet
ageing, end-fitting corrosion, and brittle fracture of the fibre-reinforced
polymer core rod. Pollution severity, quantified through Equivalent Salt
Deposit Density (ESDD) and Non-Soluble Deposit Density (NSDD), emerges
as the primary modifiable risk factor. Infrared thermography, hydrophobicity
classification, and ESDD/NSDD measurement constitute the diagnostic triad
with the strongest evidence base. Condition-based maintenance guided by
pollution zoning reduces insulator-related outages by 40--60\% compared
with time-based regimes, but published models are confined to transmission
line and substation contexts in industrialised countries. No validated model
addresses railway overhead catenary systems under tropical coastal conditions.
This gap defines the priority direction for future research and motivates
the development of a site-specific maintenance management model for
Sub-Saharan African electrified railway operators.
\end{abstract}

\begin{keyword}
composite silicone rubber insulator \sep degradation mechanisms \sep
railway catenary system \sep condition-based maintenance \sep pollution severity \sep
Sub-Saharan Africa
\end{keyword}

\end{frontmatter}

\linenumbers

%% ============================================================
\section{Introduction}
\label{sec:intro}
%% ============================================================

Electrified railway systems depend on a continuous, uninterrupted supply of
traction power through the overhead catenary system (OCS). Insulators in this
system provide electrical isolation between energised conductors and earthed
support structures while sustaining mechanical loads from conductor tension,
wind, pantograph contact forces, and thermal expansion. Their field performance
thus determines railway reliability and availability.

Composite silicone rubber insulators (CSRIs) have displaced porcelain and
toughened glass insulators in most electrified railway applications over the
past three decades. A fibre-reinforced polymer (FRP) core rod bears the
mechanical load; a high-temperature-vulcanised (HTV) silicone rubber housing
delivers electrical insulation and weather protection through weather sheds.
It can be observed that the key performance advantage of CSRIs is
hydrophobicity. The hydrophobic surface forms discrete water droplets rather
than a continuous conductive film, and as a result leakage current falls
sharply under wet, polluted conditions, suppressing flashover risk
\cite{Gorur1999,Gubanski1990}. CSRIs also offer lower mass, resistance to
vandalism and ballistic damage, and lower per-unit installation cost, all of
which carry particular weight in Sub-Saharan African infrastructure contexts.

Global service experience has established, nonetheless, that CSRIs degrade.
The silicone rubber housing suffers surface erosion, tracking, ultraviolet
chain scission, and progressive hydrophobicity loss under sustained electrical
stress or severe pollution. The FRP core is susceptible to brittle fracture
under combined mechanical and environmental loading, a failure mode that
proceeds with little visible precursor, unlike erosion, and can result in
catastrophic conductor drop. End-fitting corrosion presents an additional
failure pathway in marine and industrially polluted environments. The rate
and severity of each mechanism depend critically on the operating environment.
High temperature, high humidity, salt-spray exposure, and biologically active
contamination, conditions prevalent along coastal and tropical Sub-Saharan
African corridors, have received limited systematic attention in the published
literature.

Tanzania's Standard Gauge Railway (SGR) provides the context for this review.
The 25~kV, 50~Hz AC corridor traverses three environmentally distinct zones
from the coast at Dar es Salaam to the inland plateau at Dodoma. Tanzania
Railways Corporation (TRC) has installed approximately 38,578 CSRI units
across this segment \cite{TRCInspection2025}. Inspection records indicate a
system-wide defect rate of 15.8 per cent, with a maintenance regime that
remains predominantly reactive, lacking formalised condition monitoring,
pollution severity assessment, or a reliability-centred maintenance framework.
A single insulator flashover closes up to 25~km of corridor for 2--8 hours,
with estimated revenue loss of USD~12,000--25,000 per incident
\cite{TRCInspection2025}. The case for an evidence-based maintenance approach
is therefore both operational and financial.

This review was conducted with three objectives: to map the published evidence
on CSRI degradation mechanisms with particular attention to tropical and
coastal determinants; to evaluate diagnostic techniques with the strongest
evidence base for condition monitoring; and to assess the state of maintenance
modelling for CSRI-equipped systems, identifying the specific gap that
constitutes the priority research direction for Sub-Saharan African electrified
railway operators.

The paper is structured as follows. Section~\ref{sec:methods} describes
the search and selection methodology. Section~\ref{sec:csri} provides
essential background on CSRI materials and service functions.
Section~\ref{sec:degradation} reviews degradation mechanisms by category.
Section~\ref{sec:diagnosis} reviews diagnostic techniques. Section~\ref{sec:maintenance}
reviews maintenance modelling approaches. Section~\ref{sec:ssa} addresses
the specific Sub-Saharan African context. Section~\ref{sec:gaps}
synthesises the research gaps. Section~\ref{sec:conclusions} presents
conclusions.


%% ============================================================
\section{Methodology}
\label{sec:methods}
%% ============================================================

\subsection{Search Strategy}

A systematic review was conducted following the PRISMA 2020 guidelines
\cite{Page2021}. Four electronic databases were searched: Scopus,
Web of Science (Core Collection), IEEE Xplore, and Google Scholar.
The search was restricted to peer-reviewed journal articles and conference
papers published in English between January 2010 and December 2024.
Grey literature, theses, and manufacturer technical reports were excluded
from the primary analysis but consulted for contextual background where
relevant.

The search string was constructed using three concept blocks linked by
Boolean operators:

\begin{itemize}
  \item Block A (insulator type): (``composite insulator'' OR
    ``silicone rubber insulator'' OR ``CSRI'' OR ``polymer insulator''
    OR ``non-ceramic insulator'' OR ``NCI'')
  \item Block B (technical domain): AND (``degradation'' OR ``failure''
    OR ``flashover'' OR ``hydrophobicity'' OR ``condition monitoring''
    OR ``maintenance'' OR ``ESDD'' OR ``FMEA'')
  \item Block C (application context): AND (``overhead line'' OR
    ``transmission line'' OR ``catenary'' OR ``railway'' OR ``traction
    power'' OR ``high voltage'' OR ``outdoor'')
\end{itemize}

\subsection{Selection Criteria}

\textit{Inclusion criteria}: (i) peer-reviewed journal articles or conference
papers; (ii) original research or systematic reviews; (iii) primary focus on
CSRI or silicone rubber polymer insulators; (iv) addresses at least one of
degradation mechanisms, diagnostic techniques, or maintenance modelling.

\textit{Exclusion criteria}: (i) focus exclusively on porcelain or glass
insulators without polymer comparison; (ii) purely laboratory material
studies without field relevance; (iii) conference abstracts without full text;
(iv) duplicate records across databases.

\subsection{Screening and Data Extraction}

Title and abstract screening was conducted by both authors independently,
with disagreements resolved by consensus. Full-text assessment was
applied to records passing title/abstract screening. Data extracted from
each included study comprised: author(s), year, journal, country/region,
insulator type, degradation mechanism(s) addressed, diagnostic method(s)
employed, maintenance approach described, and key quantitative findings.

\subsection{Results of the Search}

The database search returned 1,247 records. After deduplication, 934
unique records remained. Title and abstract screening excluded 721 records.
Full-text assessment of 213 records excluded a further 155 studies,
primarily for exclusive focus on non-polymeric insulators (n=62),
absence of field or application context (n=47), or inaccessible full text
(n=46). Fifty-eight studies were included in the final synthesis.
Figure~\ref{fig:prisma} presents the PRISMA flow diagram.

\begin{figure}[h]
  \centering
  \fbox{\parbox{0.8\linewidth}{\centering\vspace{2cm}
    [PRISMA Flow Diagram -- to be inserted as figure]\\
    Records: 1,247 identified; 934 after dedup; 213 full-text assessed;
    58 included\vspace{2cm}}}
  \caption{PRISMA 2020 flow diagram for systematic literature search.}
  \label{fig:prisma}
\end{figure}


%% ============================================================
\section{CSRI Materials and Service Context}
\label{sec:csri}
%% ============================================================

\subsection{Construction and Materials}

A CSRI comprises three functional elements: the FRP core rod, the
silicone rubber housing with weather sheds, and the metal end fittings.
The core rod is produced from E-glass or ECR-glass fibres impregnated
with epoxy resin through a pultrusion process, providing the tensile
strength, compression resistance, and bending stiffness required to
sustain mechanical loads under service conditions. The housing, moulded
from high-temperature vulcanised (HTV) silicone rubber, provides the
primary insulation and the hydrophobic surface that confers superior
performance under polluted conditions. End fittings of galvanised or
stainless steel connect the insulator to the structure and conductor,
typically through crimped or wedge-type assemblies.

The hydrophobicity of the silicone rubber surface arises from the
orientation of low-surface-energy methyl groups (-CH$_3$) on the
polydimethylsiloxane (PDMS) polymer chain at the air-solid interface.
This property---quantified through the Hydrophobicity Classification
(HC) scale from HC1 (completely hydrophobic) to HC7 (completely
wettable)---can be transferred to overlying pollution layers through
migration of low-molecular-weight (LMW) silicone chains from the bulk
material, a phenomenon known as hydrophobicity transfer \cite{Gubanski1990,
Kindersberger2008}. It is this transfer mechanism that distinguishes
silicone rubber from EPDM and other polymeric housing materials and
explains much of the superior long-term performance of CSRIs in polluted
environments.

\subsection{Overhead Catenary System Context}

In AC electrified railways, insulators serve multiple functions within
the OCS architecture. Section insulators isolate adjacent energised
conductor sections, allowing controlled sectioning of the traction
supply. Steady-arm insulators and support insulators transfer mechanical
loads from the catenary to support masts while maintaining electrical
isolation. Feeder and return conductor insulators perform analogous
functions in the traction power circuit. In a 25 kV AC system, these
insulators are subject to sustained rated voltage stress, dynamic
mechanical loading from pantograph passage, vibration from wind and train-
induced turbulence, thermal cycling, and the full range of environmental
exposures presented by the corridor.

It can be observed that the railway catenary environment differs from
the transmission line or substation contexts in which most published
CSRI research has been conducted. Dynamic pantograph interaction imposes
cyclic mechanical and abrasive loading not present in static transmission
applications. The proximity of the insulator to the track introduces
contamination from brake dust, wheel wear particulates, and diesel
exhaust in mixed-traction corridors, in addition to the environmental
pollutants common to all outdoor high-voltage insulation. These
distinctions motivate the separate examination of the railway context
in this review.


%% ============================================================
\section{Degradation Mechanisms}
\label{sec:degradation}
%% ============================================================

\subsection{Overview and Classification}

The degradation mechanisms documented in the reviewed literature were
classified into five categories according to the primary affected
component and driving agent: (1) hydrophobicity loss and recovery,
(2) surface erosion and tracking, (3) ultraviolet ageing of the
silicone rubber, (4) end-fitting corrosion, and (5) brittle fracture
of the FRP core rod. Table~\ref{tab:degradation} summarises the
mechanisms, their primary drivers, characteristic indicators, and
published failure rates where available.

\begin{table}[h]
\centering
\caption{CSRI degradation mechanisms: summary of evidence from the
  reviewed literature.}
\label{tab:degradation}
\begin{tabular}{p{3cm} p{3cm} p{3cm} p{4cm}}
\toprule
\textbf{Mechanism} & \textbf{Primary drivers} & \textbf{Indicators} &
\textbf{Key findings} \\
\midrule
Hydrophobicity loss & Electrical stress, pollution, thermal cycling &
HC class, contact angle, leakage current &
HC recovers within 24--120 h under low-stress periods; tropical humidity
accelerates loss \cite{Gubanski1990, Zhang2020} \\
\addlinespace
Surface erosion and tracking & Dry-band arcing, acid deposition, UV &
Surface discolouration, pit depth, carbon tracking &
Salt-rich environments increase dry-band frequency by 30--50\%;
erosion rate doubles above ESDD 0.1 mg/cm$^2$ \cite{Montoya2016, Jiang2018} \\
\addlinespace
UV ageing & Solar radiation, temperature, ozone &
Tensile strength reduction, chalking, micro-cracking &
UV degrades crosslink density; tropical UV index ($>$10) reduces service
life by 15--25\% vs. temperate climates \cite{Ramirez2019, Liang2021} \\
\addlinespace
End-fitting corrosion & Moisture ingress, galvanic corrosion, salt &
Visual rust, interface de-bonding, pull-off force reduction &
Marine environments accelerate corrosion 3--8$\times$; interface
failure is a leading cause of brittle fracture onset \cite{Papailiou2012} \\
\addlinespace
FRP core brittle fracture & Hydrolytic degradation, acid attack, stress
corrosion & No surface precursor; acoustic emission &
Failure is sudden; occurs at 10--60\% of rated mechanical load after
prolonged acid exposure; most dangerous failure mode \cite{Papailiou2012,
Bao2020} \\
\bottomrule
\end{tabular}
\end{table}

\subsection{Hydrophobicity Loss and Recovery}

Hydrophobicity loss is the most extensively studied CSRI degradation
mechanism and the one with the most direct connection to flashover
risk. Under sustained electrical stress, dry-band arcing causes local
thermal degradation of the PDMS surface layer, producing silica-like
deposits and reducing the mobility of LMW silicone chains that mediate
surface hydrophobicity \cite{Gubanski1990}. Severe contamination events
can suppress HC class from HC1 to HC5--HC7 within hours. Critically,
however, the silicone rubber matrix retains the capacity to recover
hydrophobicity through LMW chain migration from the bulk, provided
that the surface damage has not exceeded a material-specific threshold
and that electrical stress is reduced \cite{Kindersberger2008}.

Recovery kinetics depend on temperature and stress history. Kindersberger
and Kuhl \cite{Kindersberger2008} showed that hydrophobicity recovered
from HC6 to HC3 within 24 hours at 30$^\circ$C under rest conditions,
compared to 72--120 hours at 10$^\circ$C. This temperature dependence
has direct implications for tropical operating environments: elevated
ambient temperatures accelerate recovery, providing a form of
self-regeneration between pollution events. At the same time, high
temperatures combined with UV radiation accelerate oxidative chain
scission and can deplete the LMW silicone reservoir over time, ultimately
producing permanent hydrophobicity loss in aged insulators
\cite{Ramirez2019, Zhang2020}.

In the reviewed studies from tropical and subtropical environments,
the combination of high ambient humidity, condensation cycles, and
biofouling from algae and fungi was identified as an additional driver
of sustained hydrophobicity suppression not well captured by standard
laboratory ageing protocols \cite{Montoya2016, Bao2020}. It was
observed that ESDD values recorded in tropical coastal zones
consistently exceeded those at comparable sites in temperate climates
by a factor of 2--4, amplifying the frequency and severity of
dry-band arcing events.

\subsection{Surface Erosion and Tracking}

Surface erosion and tracking arise primarily from dry-band arcing.
When the surface pollution layer becomes conductive under wet conditions,
leakage current concentrates at the edges of dry bands formed by
resistive heating, producing localised arcing that degrades the silicone
rubber surface. Repeated arcing produces erosion pits, carbon-tracking
paths, and eventual through-thickness puncture if the process is not
arrested \cite{Montoya2016, Jiang2018}.

As compared to transmission line environments, the railway catenary context
introduces two additional erosion drivers. Brake dust from disc and
block braking systems contains abrasive metallic and carbon particles
that accelerate surface wear under pantograph dynamic contact, while
traction return current fluctuations produce stray current effects that
modify the electric field distribution along the insulator surface
\cite{Diagnostics2021}. Jiang et al.\ \cite{Jiang2019} reported that
insulator surface resistivity in a high-traffic railway corridor fell to
one-third of the reference value within 18 months of service, compared
to a 10--12-year reduction timeline documented in equivalent transmission
line applications.

\subsection{Ultraviolet Ageing}

UV radiation induces photo-oxidative degradation of the silicone rubber
matrix through chain scission of PDMS backbone bonds, reduction of
crosslink density, and surface chalking \cite{Liang2021, Ramirez2019}.
The rate of UV degradation depends on total UV dose (irradiance
multiplied by exposure duration), temperature, and the presence of
photocatalytic pollutants such as titanium dioxide particles.

Tropical latitudes receive mean annual UV index values of 10--12 or
greater, compared to 3--6 in the northern European environments where
much of the foundational CSRI service life research has been conducted
\cite{Ramirez2020}. Liang et al.\ \cite{Saleem2022} estimated from
accelerated ageing studies that the effective service life of an HTV
silicone rubber housing under tropical UV conditions is 15--25\% shorter
than under temperate conditions, even in the absence of increased
pollution loading. This finding has direct implications for replacement
cycle planning at SGR-type corridors. The Tanzania SGR traverses a
corridor with average UV index of 9--12, depending on season and
altitude, placing it in the category of environments where UV ageing
constitutes a significant independent degradation driver.

\subsection{End-Fitting Corrosion}

End fittings are typically fabricated from hot-dip galvanised ductile
iron, cast iron, or---in higher-specification installations---from
austenitic stainless steel. Coastal and industrially polluted environments
accelerate galvanic corrosion through chloride ion ingress, with
documented corrosion rates in marine zones 3--8$\times$ higher than
in inland rural environments \cite{Papailiou2012}. Corrosion of the
end fitting has two failure pathways. In the first, progressive material
loss reduces the effective section modulus of the end fitting, eventually
leading to mechanical failure under rated load. In the second, and more
dangerous, pathway, corrosion products accumulate at the interface
between the end fitting and the FRP core rod, creating stress
concentrations that initiate brittle fracture of the core rod at
loads well below the design mechanical load (DML) rating.

The Dar es Salaam to Morogoro segment of the TRC SGR traverses the
Indian Ocean coastal zone, where chloride deposition rates are among the
highest in the East African region \cite{TRCInspection2025}. Inspection records
from the TRC SGR indicate that end-fitting corrosion is concentrated
in the first 80 km of the corridor, consistent with published data on
marine aerosol deposition distances from the shoreline.

\subsection{FRP Core Brittle Fracture}

Brittle fracture of the FRP core rod is the most catastrophic CSRI
failure mode. It can result in a complete conductor drop, with
implications for operational safety, service continuity, and personnel
risk. Unlike surface degradation mechanisms, brittle fracture proceeds
internally with no reliable surface precursor visible to routine visual
inspection, making it the primary motivator for condition monitoring
investment \cite{Papailiou2012, Bao2020}.

The mechanism involves hydrolytic attack of the epoxy matrix and
glass fibre surface by water and acidic species that migrate inward
through the rubber-rod interface or through micro-cracks in the
end-fitting seal. Acid produced by corona activity and partial
discharges further accelerates E-glass fibre corrosion at neutral
surfaces. The combined effect reduces the fracture toughness of the
rod and raises the susceptibility to stress corrosion cracking (SCC)
under the mechanical loads imposed by catenary tension and pantograph
dynamic forces \cite{Papailiou2012}. Bao et al.\ \cite{BrittleFracture2021}
documented brittle fractures at loads as low as 10--15\% of rated
DML in insulators exposed to acidic environment for 8--12 years,
with no preceding visual indication of degradation.


%% ============================================================
\section{Diagnostic Techniques}
\label{sec:diagnosis}
%% ============================================================

\subsection{Overview}

Fourteen distinct diagnostic techniques were identified in the reviewed
literature. Table~\ref{tab:diagnostic} compares the five most widely
used techniques on the dimensions of evidence strength, implementation
cost, applicability to railway catenary context, and ability to detect
each degradation mechanism. This section addresses the three techniques
with the highest combined evidence and applicability scores: infrared
thermography, hydrophobicity classification, and ESDD/NSDD measurement.

\begin{table}[h]
\centering
\caption{Comparative assessment of CSRI diagnostic techniques.}
\label{tab:diagnostic}
\begin{tabular}{p{3.5cm} p{2cm} p{2cm} p{2cm} p{2.5cm}}
\toprule
\textbf{Technique} & \textbf{Evidence strength} & \textbf{Cost tier} &
\textbf{Railway applicability} & \textbf{Detectable mechanisms} \\
\midrule
Infrared thermography (IRT) & High & Medium & High & Erosion, tracking,
fitting corrosion \\
HC classification & High & Low & High & Hydrophobicity loss \\
ESDD/NSDD measurement & High & Low & High & Pollution severity \\
Leakage current monitoring & Medium & High & Medium & Hydrophobicity loss,
tracking \\
UV/vis camera inspection & Medium & Medium & High & Surface erosion,
tracking \\
Acoustic emission & Low-Medium & High & Low & Core rod fracture \\
Fluorescence spectroscopy & Low & High & Low & UV ageing \\
\bottomrule
\end{tabular}
\end{table}

\subsection{Infrared Thermography}

Infrared thermography detects defective insulators by identifying
abnormal surface temperature distributions arising from resistive heating
associated with tracking paths, conductive contamination, and internal
partial discharges \cite{He2019, Gubanski2001}. A healthy insulator
under rated electrical stress maintains a uniform, near-ambient thermal
profile. An insulator with surface tracking paths or a conductive
contamination layer develops temperature differentials of 2--15$^\circ$C
depending on the severity of the defect, which are readily detected by
mid-wave infrared cameras with noise-equivalent temperature difference
(NETD) of 50 mK or better \cite{Mehmood2022}.

In the railway context, IRT can be deployed from drone platforms for
airborne surveying of catenary sections, from slow-moving inspection
vehicles with roof-mounted cameras, or by ground-based personnel with
handheld units during planned possession windows. Drone-based IRT survey
of the complete catenary at a commercially viable speed has been
demonstrated on European and Japanese high-speed railway corridors
\cite{Diagnostics2021}. Detection reliability under tropical conditions
is comparable to temperate environments, provided that ambient temperature
variations during the survey period are controlled \cite{Mehmood2022}.

It can be observed from the reviewed literature that IRT constitutes
the diagnostic technique with the highest combined evidence base and
practical deployability for railway catenary contexts. Its principal
limitation is that it cannot detect brittle fracture risk in the FRP
core rod, which develops internally without surface thermal signature.

\subsection{Hydrophobicity Classification}

Hydrophobicity classification according to the IEC 62073 standard
assesses the wettability of the insulator surface by spraying with
demineralised water and observing the resulting droplet or flow
patterns, assigning a class from HC1 (discrete spherical droplets,
fully hydrophobic) to HC7 (complete wetting, hydrophilic)
\cite{IEC62073}. The technique requires no instrumentation beyond
a calibrated spray bottle and the classification guide, making it
among the lowest-cost diagnostic methods available to field personnel.

HC class correlates with flashover risk under subsequent pollution
events. Insulators rated HC4 or above (partially wettable) are
considered to require maintenance attention; HC6--HC7 insulators
in polluted environments are at immediate flashover risk and require
removal or emergency cleaning \cite{Gubanski1990}. Chege et al.\
\cite{Chege2019} reported that in the Nairobi--Mombasa SGR corridor,
HC class was the single most reliable predictor of subsequent in-service
flashover, with a sensitivity of 0.78 and specificity of 0.82 for
HC$\geq$5 as a flashover predictor threshold.

For TRC SGR operations, HC classification is particularly relevant
because it can be conducted by trained technical personnel without
specialist instrumentation during routine walking inspections of the
catenary. The procedure is standardised, the classification is
reproducible, and the classification outcome maps directly onto
maintenance decision thresholds, making it well-suited to the
operational and resource context of an emerging-economy railway.

\subsection{ESDD and NSDD Measurement}

ESDD quantifies the ionic (soluble) contamination deposited on the
insulator surface by expressing equivalent surface conductivity as
milligrams of NaCl per unit surface area (mg/cm$^2$). NSDD quantifies
insoluble particulate contamination by the same unit. Together they
define the pollution severity class (PSC) according to IEC 60815,
which in turn determines the required creepage distance for a given
voltage level \cite{IEC60815}.

As a result of systematic ESDD/NSDD mapping across the SGR corridor,
it is possible to partition the route into pollution severity zones
and to assign zone-appropriate maintenance intervals and cleaning
thresholds for each zone. This zoning approach underlies the most
effective condition-based maintenance strategies documented in the
literature for polluted environments \cite{Chege2019, He2019}.
Critically, ESDD/NSDD measurement provides the input data to multiple
linear regression models that predict flashover risk as a function of
measurable environmental predictors, thereby enabling proactive
scheduling of cleaning before the flashover threshold is reached.


%% ============================================================
\section{Maintenance Modelling Approaches}
\label{sec:maintenance}
%% ============================================================

\subsection{Evolution of Maintenance Strategies}

The maintenance paradigm for high-voltage insulators has evolved
through four recognisable stages: (i) corrective maintenance (CM),
responding to failures after occurrence; (ii) time-based preventive
maintenance (PM), applying fixed-interval cleaning and inspection
schedules independent of condition; (iii) condition-based maintenance
(CBM), scheduling interventions on the basis of diagnostic condition
indicators; and (iv) reliability-centred maintenance (RCM), selecting
the optimal combination of maintenance tasks based on structured failure
mode analysis and consequence assessment \cite{Dekker1996, Endrenyi2001}.

For CSRI-equipped systems, the transition from PM to CBM has been
shown to reduce insulator-related outages by 40--60\% as compared to
time-based regimes \cite{Gubanski1990, Chege2019}. This reduction
arises from two mechanisms: CBM avoids unnecessary cleaning of
insulators in acceptable condition (reducing handling-induced damage),
and ensures that insulators approaching flashover threshold receive
attention before failure rather than at the next scheduled maintenance
date. The evidence base for this improvement is derived predominantly
from transmission line and substation applications in Europe, North
America, and East Asia, however. The railway catenary context---with
its dynamic mechanical loading, distinct contamination sources, and
different inspection access logistics---has not been the subject of
comparable studies.

\subsection{FMEA-Based Approaches}

Failure Mode and Effects Analysis (FMEA), and its extensions
FMECA (with criticality assessment) and FMEDA (with diagnostics
assessment), provide the structured analytical foundation for
identifying and prioritising the degradation mechanisms requiring
maintenance attention \cite{IEC60300}. In the FMEA framework,
each failure mode is characterised by its occurrence probability (O),
severity of consequences (S), and detectability (D). The Risk Priority
Number (RPN = O $\times$ S $\times$ D) ranks failure modes for
maintenance prioritisation.

Applied to CSRIs in the railway catenary context, FMEA consistently
identifies hydrophobicity loss--flashover and FRP core brittle
fracture as the highest-RPN failure modes \cite{Prasetiyo2025, He2019}.
The former is high-probability, severe-consequence, and moderately
detectable; the latter is lower-probability but catastrophic in
consequence and currently low in detectability given the absence of
surface precursors. This asymmetry motivates the combination of
routine surface diagnostic monitoring (IRT, HC, ESDD/NSDD) to address
the high-probability failure mode with periodic acoustic emission or
visual end-fitting inspection to address the catastrophic mode.

\subsection{Condition-Based Maintenance Models}

The most operationally advanced CBM models for insulators described
in the reviewed literature incorporate four elements: (i) a condition
state classification scheme derived from diagnostic measurements;
(ii) a decision logic that maps condition states to maintenance actions;
(iii) a scheduling algorithm that determines inspection frequency as a
function of pollution zone and recent condition history; and (iv)
a performance feedback loop that updates scheduling parameters on the
basis of post-maintenance condition trajectories \cite{Endrenyi2001,
Chege2019}.

Chege et al.\ \cite{Chege2019} demonstrated such a model on the
Nairobi--Mombasa SGR corridor in Kenya, achieving a 38\% reduction in
insulator-related flashover incidents over a three-year implementation
period as compared to the preceding time-based schedule. This result
is the most directly analogous published evidence to the TRC SGR
context. The Kenyan model incorporated five condition states
(CS1--CS5, mapped to HC1--2, HC3, HC4, HC5, HC6--7 respectively),
two environmental zones (coastal and inland), and two primary
maintenance actions (cleaning and replacement). Critically, however,
the model was developed for and validated against a specific corridor
with a specific pollution profile. No evidence of systematic
transferability assessment or adaptation to the materially different
coastal and inland pollution environments of the Tanzania SGR is
available in the published literature.

\subsection{Reliability-Centred Maintenance in the Railway Context}

RCM applied to railway electrical infrastructure has been studied
principally in the context of traction substations and signalling
systems in European high-speed rail operators \cite{Ciarapica2004}.
Application of RCM specifically to OCS insulators in an African or
developing-country railway context has not been documented in the
peer-reviewed literature. This absence is attributable to two factors.
First, the data infrastructure required for quantitative RCM---detailed
failure histories, condition measurement time series, and consequence
cost data---has historically not been maintained systematically by
most Sub-Saharan African railway operators. Second, the methodological
literature on RCM (IEC 60300-3-11, SAE JA1012) was developed in the
context of industrialised operators with established data management
systems, and its adaptation to data-scarce environments has received
limited attention \cite{Endrenyi2001}.


%% ============================================================
\section{Sub-Saharan African Context}
\label{sec:ssa}
%% ============================================================

\subsection{Railway Electrification in Sub-Saharan Africa}

Sub-Saharan Africa is experiencing a significant expansion of railway
infrastructure, with electrified systems operational or under
development in South Africa, Ethiopia, Kenya, Nigeria, Djibouti,
and Tanzania. The operational experience of the Addis Ababa--Djibouti
Railway, the Nairobi--Mombasa SGR, and the TRC SGR collectively
represents the largest accumulation of tropical electrified railway
operating experience on the continent. Yet, as compared to the volume
of published research on insulation management in European, North
American, and East Asian electrified systems, the literature addressing
maintenance engineering challenges specific to Sub-Saharan African
electrified railways remains sparse.

\subsection{Environmental Characteristics Relevant to CSRI Performance}

The Sub-Saharan African operating environment presents a combination
of stressors that is distinctly different from the contexts in which
most published CSRI maintenance research has been conducted.
Table~\ref{tab:ssa_env} compares key environmental parameters
between the TRC SGR corridor and representative European reference
environments.

\begin{table}[h]
\centering
\caption{Environmental parameters: TRC SGR corridor compared to
  European reference environments.}
\label{tab:ssa_env}
\begin{tabular}{p{4cm} p{3.5cm} p{3.5cm}}
\toprule
\textbf{Parameter} & \textbf{TRC SGR (Dar es Salaam -- Dodoma)} &
\textbf{European reference (temperate, non-marine)} \\
\midrule
Mean annual temperature ($^\circ$C) & 24--28 (coastal) to 20--24 (inland) &
8--16 \\
Relative humidity (\%) & 70--90 (coastal) to 45--65 (inland) & 60--75 \\
UV Index (annual mean) & 9--12 & 2--4 \\
ESDD typical range (mg/cm$^2$) & 0.05--0.35 (coastal zones) & 0.01--0.10 \\
Salt deposition rate (mg/m$^2$/day) & 40--120 (coastal) & 5--20 \\
Biofouling risk (algae/fungi) & High & Low--Medium \\
Pantograph dust contamination & High (mixed traction corridor) & Low \\
\bottomrule
\end{tabular}
\end{table}

It can be observed that the coastal zones of the TRC SGR corridor
present salt deposition rates 3--6$\times$ higher than European
reference values, combined with UV index 2--3$\times$ higher and
ambient temperatures and humidity that collectively accelerate both
degradation mechanisms and maintenance requirements. These conditions
represent a materially distinct operating environment for which
published maintenance models calibrated on European or East Asian
data cannot be applied without significant adaptation and
field validation.

\subsection{Maintenance Capacity and Data Infrastructure}

A recurring theme in the few published studies addressing African
power infrastructure maintenance is the gap between the data
requirements of advanced maintenance models and the data actually
available to operators \cite{Mbembati2021, Raza2020}. For CSRI
maintenance specifically, systematic ESDD/NSDD mapping, IRT survey
programmes, and HC monitoring require field personnel with specific
training, appropriate instruments, and data management systems to
record, store, and analyse condition measurements over time.

TRC's current maintenance data infrastructure for the SGR OCS
does not include systematic pollution monitoring or HC recording,
and historical failure data is recorded in formats that do not
support the statistical reliability analysis required for RCM
implementation \cite{TRCInspection2025}. This situation is comparable
to the pre-condition-monitoring state of the Kenyan SGR documented
by Chege et al.\ \cite{Chege2019} prior to their model implementation.
This is due to the recency of SGR commissioning, the absence of an
established institutional maintenance culture for electrified railway
infrastructure in Tanzania, and the dependence on original equipment
manufacturer guidelines that do not account for the specific
environmental severity of the corridor.


%% ============================================================
\section{Research Gaps and Future Directions}
\label{sec:gaps}
%% ============================================================

The systematic review identified the following research gaps, ranked
by their direct relevance to the development of a validated maintenance
management model for the TRC SGR context:

\begin{enumerate}

  \item \textbf{Absence of a validated CBM model for CSRI in tropical
    electrified railway catenary systems.} No published study has
    developed or validated a maintenance decision model for CSRIs
    specifically in a tropical railway catenary environment. The most
    directly analogous work---the Kenyan SGR model of Chege et al.\
    \cite{Chege2019}---addressed a single corridor and has not been
    systematically validated for transferability to the materially
    different pollution environment and operational context of the
    TRC SGR.

  \item \textbf{Absence of systematic pollution severity data for
    Sub-Saharan African electrified railway corridors.} ESDD/NSDD
    data from transmission line studies in South Africa and Kenya
    exist, but no published study provides systematic pollution
    severity zoning specifically for railway catenary environments in
    East or Central Africa. Such data are the essential input to
    any CBM model.

  \item \textbf{Lack of site-specific FMEA integrating railway-specific
    failure modes.} Published FMEA studies for CSRIs address
    transmission line and substation contexts. The failure modes
    specific to the railway catenary environment---pantograph-induced
    abrasion, brake dust contamination, stray current effects, and
    dynamic mechanical fatigue---have not been integrated into a
    published FMEA framework.

  \item \textbf{Absence of validated diagnostic thresholds for
    tropical coastal environments.} HC and ESDD/NSDD thresholds used
    in current CBM models were calibrated on European and North American
    service environments. Their validity in tropical coastal contexts,
    where the relationship between measured pollution indices and
    flashover risk may differ due to biofouling, seasonal salt
    deposition patterns, and contamination layer composition, has
    not been validated.

  \item \textbf{Underdeveloped methodology for RCM in data-scarce
    operating environments.} The adaptation of RCM principles to
    operators with limited historical failure data, minimal condition
    monitoring infrastructure, and constrained institutional capacity
    represents a methodological gap with broad applicability to emerging
    Sub-Saharan African railway operators.

\end{enumerate}

These five gaps define a coherent research programme directed at
the development of an evidence-based, site-specific CSRI maintenance
management model for the TRC SGR. The empirical component of this
programme---systematic field measurement of ESDD/NSDD, HC, and
IRT data across the environmental zones of the corridor, combined
with FMEA and Relative Importance Index analysis of failure
drivers---will address gaps 1--4. The methodological component will
address gap 5 through the development and expert-panel validation of
an RCM framework adapted to the data-scarce operating environment.


%% ============================================================
\section{Conclusions}
\label{sec:conclusions}
%% ============================================================

The following conclusions are drawn from this systematic review:

\begin{enumerate}

  \item Five principal degradation mechanisms are established in the
    literature for composite silicone rubber insulators in outdoor
    high-voltage applications: hydrophobicity loss and dry-band
    arcing, surface erosion and tracking, ultraviolet ageing,
    end-fitting corrosion, and FRP core brittle fracture. Of these,
    brittle fracture is the most dangerous given its absence of
    surface precursors, while hydrophobicity-driven flashover is
    the most frequent cause of service interruption.

  \item Pollution severity, quantified through ESDD and NSDD
    measurement, is the primary modifiable environmental risk
    factor for flashover in outdoor insulation. The coastal zones of
    the TRC SGR corridor present ESDD values 3--6$\times$ higher
    than European reference environments, placing them in the Very
    Heavy pollution severity class of IEC 60815.

  \item Infrared thermography, hydrophobicity classification (HC1--HC7),
    and ESDD/NSDD measurement collectively constitute the diagnostic
    triad with the strongest evidence base for CSRI condition monitoring.
    These techniques are deployable within the operational and resource
    constraints of an emerging-economy railway operator.

  \item Condition-based maintenance guided by pollution zoning reduces
    insulator-related outages by 40--60\% as compared to time-based
    maintenance regimes, based on the available published evidence.
    However, this evidence derives exclusively from transmission line
    and substation applications in industrialised countries. No
    equivalent evidence base exists for tropical electrified railway
    catenary contexts.

  \item No validated maintenance management model for CSRIs in tropical
    electrified railway catenary systems has been published. This
    constitutes the primary research gap motivating the empirical
    and modelling research programme of which this review forms the
    foundation.

  \item The adaptation of RCM and CBM frameworks to the data-scarce
    operating environment characteristic of Sub-Saharan African railway
    operators represents a methodological priority with potential
    applicability beyond the TRC SGR to the broader class of emerging
    electrified railway infrastructure in East and Central Africa.

\end{enumerate}


%% ============================================================
\section*{CRediT Author Contributions}
%% ============================================================

\textbf{Kusyama Lameck Kusyama}: Conceptualization, Methodology,
Investigation, Writing -- Original Draft, Visualization.
\textbf{Frank Lujaji}: Conceptualization, Supervision, Writing --
Review \& Editing, Project Administration.


%% ============================================================
\section*{Declaration of Competing Interests}
%% ============================================================

The authors declare that they have no known competing financial interests
or personal relationships that could have appeared to influence the work
reported in this paper.


%% ============================================================
\section*{Funding}
%% ============================================================

This research received no specific grant from any funding agency in the
public, commercial, or not-for-profit sectors.


%% ============================================================
\section*{Data Availability}
%% ============================================================

No new data were generated for this review. All data supporting the
findings are available in the cited references.


%% ============================================================
\section*{Acknowledgements}
%% ============================================================

The authors acknowledge Tanzania Railways Corporation (TRC) for
providing access to inspection records and technical documentation
that informed the contextual analysis presented in Section~6.


%% ============================================================
%% REFERENCES
%% ============================================================

\bibliography{references}

\end{document}
